% ===== §3 =====
\section{From Concepts to Persons}
\label{p27:sec:bridge}

\noindent
The step from the epistemology to the present subject is short, and the companion work took it in a sentence that the present work will unfold. A discipline is a symbolic system, and the crossing of disciplines affords the observation of a concept that several disciplines concern. A person is also the bearer of a symbolic system. Each person renders the world in a way shaped by an experience, a history, a native language, a family, a culture, a body, that no other person exactly shares, so that two people who use what appears to be the same word may render beneath it two aspects that do not coincide. The dialogue between persons is therefore a crossing of symbolic worlds in the same sense that interdisciplinary inquiry is, differing in that the worlds crossed are borne by living persons where the systems of inquiry are codified in fields, and carried with the whole weight of a life where a literature sets them down.

What the crossing affords is the same in both cases. It is observation, won through the relations among partial renderings of a matter that exceeds them all. When two persons bring their renderings of a shared matter, a memory, a fear, a hope, the meaning of a word like home or freedom or love, into relation, the matter becomes observable in a way that neither rendering held alone. The understanding they reach is not the recovery of what one of them already meant, transported intact to the other. It is a generative observation, produced in the relation between two worlds, and belonging to neither.

\begin{claim}[The relational form of understanding another]
To understand another person is not to decode a message the other has sent, recovering intact a meaning the other already possessed. It is to bring the renderings of two symbolic worlds into relation and to observe, in that relation, a matter that neither world held alone. Understanding another is a generative observation, produced between two worlds and belonging to neither.
\end{claim}

Intimacy is the domain in which this crossing is carried furthest. Two people who share a life bring their symbolic worlds into a relation more sustained, more entangled, and more consequential than any other bond affords. They render to one another the whole texture of a shared existence and not isolated concepts, and they do so over years, so that the crossing is a long history through which each world is altered by its relation to the other, and not a single act. The depth of intimacy is the depth of this crossing, and its difficulty is the difficulty of the crossing carried out where the stakes are a shared life.

This depth carries a corresponding danger, and naming it here prepares the whole of what follows. Because the two worlds in an intimate bond are so entangled, and because the partners so often use the same words, there arises a powerful illusion that they inhabit a single symbolic world, that no crossing is required, that to understand the other is simply to consult oneself. This illusion is the source of a characteristic and painful form of failure, the discovery, often late, that what one took for a shared meaning was two divergent renderings all along, and that a bond one believed transparent was crossed by an unrecognized difference of worlds. The account developed here holds that the recognition of the other as the bearer of a distinct symbolic world, so easily lost in the very closeness of intimacy, is the first condition of understanding within it, and that much of what passes for the failure of love is in truth the failure of this recognition.

\subsection*{The unsymbolizable core}

There is a feature of the other's symbolic world that must be introduced here, in the background, since a later part of the paper turns upon it. The account has said that a symbolic world is never exhausted by another's observation of it, and this holds with a particular force for a person. Psychoanalysis, in the work of Lacan~\parencite{lacan1977xi,lacan1992vii}, gives this a precise shape. The subject who enters language does so at the cost of a remainder that language does not carry, a core that the signifier cannot reach, and around this core the subject's desire is organized. Lacan marks the remainder in several registers, as the object-cause of desire that no possessed object supplies, as the lack in the subject that the entry into speech installs, as a real that resists symbolization. The details of the doctrine are not required here, and the series has developed them elsewhere. What the present paper takes from it is a single structural point, that at the center of the other's symbolic world there is a core that no crossing reaches, not through a want of skill or a shortage of time, but because the core is of a kind that the signifier cannot render.

This structural point sharpens the epistemology's claim of incompleteness in the case of a person. When the paper says that the observation of the other is always incomplete, it does not mean only that there is always more of the other to observe, a further region that a longer acquaintance would reach. It means, in addition, that at the heart of the other there is a core of another kind, no further region to be reached at all, an unsymbolizable center around which the other's desire moves and which the other does not possess as a content to be conveyed. The incompleteness of understanding another is therefore structural as well as extensive, a matter of a core that observation approaches and does not close, and not only a matter of how much remains to be observed. This core, introduced here, returns as the pivot of the paper's dialectical turn, where it is shown that the persistence of the bond depends upon it.
